\documentclass[a4paper,12pt]{article}

\usepackage[T2A]{fontenc}
\usepackage[utf8]{inputenc}
\usepackage[ukrainian]{babel}
\usepackage{geometry}
\usepackage{anyfontsize}
\usepackage{setspace}
\usepackage{hyperref}

\geometry{margin=2cm}
\setstretch{1.15}

\AtBeginDocument{\fontsize{14}{17}\selectfont}

\begin{document}

\begin{titlepage}
\begin{center}
Мiнiстерство освiти i науки України\\
НАЦIОНАЛЬНИЙ УНIВЕРСИТЕТ\\
\textquotedblleft КИЄВО-МОГИЛЯНСЬКА АКАДЕМIЯ\textquotedblright \\
Кафедра мультимедiйних систем факультету iнформатики\\[7em]
\textbf{ПРОМIЖНИЙ ЗВIТ}\\[0.4em]
до курсової роботи за спецiальнiстю \\\textquotedblleft Iнженерiя програмного забезпечення\textquotedblright \\[0.6em]
\textbf{Тема: Трасування променiв у воксельному просторi}\\[2.0em]
\end{center}

\begin{flushleft}
Виконавець: студент 3 курсу факультету iнформатики, Олександр Гнутов\\
Керiвник: старший викладач Борозенний С. О.
\end{flushleft}

\vfill
\begin{center}
Київ 2026\\
31 сiчня 2026
\end{center}
\end{titlepage}

\section*{Мета та завдання роботи}
Метою курсової роботи є дослiдити можливiсть застосування трасування променiв у дискретному воксельному просторi та реалiзувати демонстрацiйний застосунок, який виконує обчислення тiней та наближеного навколишнього перекриття у воксельному представленнi сцени на GPU.

У межах поточного етапу сформульовано та деталiзовано такi завдання:
\begin{enumerate}
\item Проаналiзувати предметну область: воксельнi структури даних, методи проходження променя у решiтцi (DDA та Amanatides--Woo), а також прийоми прискорення на GPU.
\item Побудувати прототип рендерера на базi рушiя Bevy мовою Rust, що забезпечує завантаження воксельних даних та керування сценою.
\item Реалiзувати загальний воксельний об’єм сцени (World Volume) у виглядi 3D текстури на GPU та сформувати MIP-рiвнi для пришвидшення пропуску порожнього простору.
\item Реалiзувати постпроцесинговий вiдкладений (deferred) етап, що використовує данi з буфера глибини та нормалей (G-buffer) для запуску трасування у world volume.
\item Додати засоби для контролю параметрiв i порiвняння варiантiв шейдера (наївний i оптимiзований), а також виведення режимiв вiзуалiзацiї.
\end{enumerate}

\section*{Короткий опис програмної реалiзацiї}
Проєкт реалiзовано мовою Rust з використанням Bevy 0.15.3. Архiтектурно застосунок побудований як набiр плагiнiв, що iнтегруються в Bevy App та RenderApp.

У файлах (src/lib.rs) i (src/main.rs) налаштовано запуск застосунку, камеру, базовi об’єкти сцени та налагоджувальний iнтерфейс. У застосунку передбачено інтерактивнi параметри (через \texttt{egui}): вибiр режиму вiзуалiзацiї, величина зсуву (bias) для уникнення самозатiнення, перемикання варiанту шейдера (naive або optimized), керування видимiстю воксельних моделей та world volume.

\subsection*{Матерiали та буфери для вiдкладеного рендерингу}
Для воксельних моделей використовується матерiал VoxelRayTracedMaterial, який зберiгає 3D текстуру з воксельними даними. На етапi рендерингу матерiал формує додатковi вихiднi ресурси: буфер глибини та буфер нормалей, якi далi використовуються постпроцесом.

У поточнiй версiї сцена мiстить декiлька прикладiв воксельних об’єктiв, зокрема дерево та рельєф, а також окремий об’єкт (умовна зброя), для якого передбачено режим ігнорування глибини.

\subsection*{World Volume як уніфікований простiр сцени}
Ключовим компонентом реалiзацiї є World Volume, який представлено як 3D текстуру формату R8Uint. Такий формат дозволяє компактно зберiгати бiнарну iнформацiю про заповненiсть вокселя. Розмiр world volume параметризується на рiвнi VoxelPlugin i в демонстрацiйнiй конфiгурацiї встановлений як 256 на 256 на 256.

Заповнення world volume здiйснюється обчислювальним шейдером. Для кожного воксельного об’єкта сцени на CPU формується матриця перетворення координат, яка переводить координати world volume до локальних координат моделi. У RenderApp виконується проходження по колекцiї активних воксельних об’єктiв та запуск compute-проходу для внесення їх даних у загальний об’єм.

Для прискорення трасування у world volume реалiзовано генерацiю MIP-рiвнiв (у поточнiй конфiгурацiї три рiвнi). Це дає змогу у шейдерi використовувати крупнiший крок проходження по решiтцi та швидше пропускати порожнi дiлянки, наближаючись до iдеї октодерева без явної побудови дерева.

\subsection*{Постпроцесингове трасування у воксельному просторi}
Пiсля основного проходу рендерингу запускається постпроцес DeferredRayTracedNode, який отримує:
\begin{enumerate}
\item екранну текстуру з базовим кольором
\item буфер глибини
\item буфер нормалей
\item world volume як 3D текстуру
\item додатковий шум або джитер для стохастичних вибірок
\end{enumerate}

У WGSL реалiзовано два варiанти фрагментного шейдера: базовий та оптимiзований. Обидва варiанти використовують данi G-buffer для вiдновлення точки в просторi для кожного пiкселя, обчислюють напрямок променя для вибiрок (для тiней та AO) та виконують проходження по воксельнiй решiтцi.

Наївний варiант використовує фiксований MIP-рiвень, що обирається через налаштування. Оптимiзований варiант реалiзує розрiджене проходження trace\_sparse, яке стартує з бiльш грубого рiвня та динамiчно зменшує MIP при наближеннi до потенцiйного перетину, тим самим зменшуючи кiлькiсть крокiв.

Також враховано проблему самозатiнення. Для цього початкова точка променя змiщується вздовж нормалi поверхнi та додатково вздовж напрямку променя на частку розмiру вокселя.

\section*{Поточний стан виконання робiт}
Станом на 31 сiчня 2026 виконано ключовi пункти етапу прототипування:
\begin{enumerate}
\item Пiдготовлено робочий каркас застосунку на Bevy з базовою сценою, камерою та налагоджувальним iнтерфейсом.
\item Реалiзовано матерiали для воксельних моделей та формування допомiжних буферiв (глибина i нормалi).
\item Реалiзовано структуру world volume та compute-проходи для очищення, заповнення та генерацiї MIP-рiвнiв.
\item Реалiзовано deferred постпроцес, який виконує трасування у world volume та моделює тiнь i наближене ambient occlusion.
\item Додано можливiсть перемикання між двома шейдерами i декiлькома режимами вiзуалiзацiї для перевiрки коректностi.
\end{enumerate}

\section*{Промiжнi результати та спостереження}
Отримано працездатний прототип, який демонструє наступне:
\begin{enumerate}
\item Єдиний глобальний об’єм дозволяє уніфікувати трасування для множини об’єктiв без побудови iндивiдуальних прискорювальних структур.
\item MIP-рiвнi дають практичний шлях до прискорення: у типових сценах з великою кiлькiстю порожнього простору оптимiзований прохід зменшує кiлькiсть звернень до textureLoad.
\item У воксельному представленнi помiтнi артефакти самозатiнення, тому bias та комбiнований зсув (по нормалi й по променю) є обов’язковими для прийнятної якостi.
\item Стохастичнi вибiрки для AO потребують компромiсу мiж кiлькiстю семплiв i швидкiстю. У поточнiй конфiгурацiї використано невелику кiлькiсть вибiрок, що дає шум, але зберiгає продуктивнiсть.
\end{enumerate}

\section*{Проблеми та ризики}
На поточному етапi iдентифiковано такi технiчнi ризики:
\begin{enumerate}
\item Обмеження за розмiрами world volume та вартiсть пам’ятi GPU. Зi збiльшенням роздiльностi об’єму зростають витрати на зберiгання та на генерацiю MIP-рiвнiв.
\item Баланс якостi й швидкодiї. Для стабiльного FPS потрiбне коректне обмеження кiлькостi крокiв трасування та вибiрок AO.
\item Вiдтворюванiсть якостi при змiнi масштабу сцени. При iнших розмiрах вокселя потрiбна повторна калибровка bias та параметрiв освiтлення.
\end{enumerate}

\section*{План робiт на наступний етап}
До наступного контрольного етапу планується:
\begin{enumerate}
\item Реалізувати додаткові оптимізації та покращення (однобітний воксельний об'єм, тіньове акне при русі об'єктів, додаткові джерела світла).
\item Оформити роздiл 2 текстової частини з описом DDA, Amanatides--Woo та логiки розрiдженого обходу на основi MIP-рiвнiв.
\item Оформити роздiл 3 з описом архiтектури застосунку: модулi Rust, граф рендерингу Bevy, формування world volume та взаємодiя з WGSL.
\item Пiдготувати набiр вiзуальних результатiв для демонстрацiї: режими нормалей, глибини, напрямку променя, а також порiвняння наївного та оптимiзованого варiанту.
\item Додати простi вимiрювання продуктивностi для рiзних розмiрiв world volume та рiзних режимiв трасування, щоб сформувати висновки щодо ефективностi оптимiзацiй.
\end{enumerate}

\section*{Висновок}
На момент промiжного контролю реалiзовано працюючий прототип трасування променiв у воксельному просторi на GPU з використанням загального воксельного об’єму та двох варiантiв шейдерної реалiзацiї. Архiтектура на базi Bevy дозволяє інтегрувати compute-етапи для world volume та фрагментний deferred-прохід у рендер-граф, а налагоджувальний iнтерфейс дає змогу оперативно перевiряти коректнiсть та параметри якостi.

\begin{thebibliography}{9}
\bibitem{wiki_voxel} Воксель. Вiкiпедiя. \url{https://uk.wikipedia.org/wiki/%D0%92%D0%BE%D0%BA%D1%81%D0%B5%D0%BB%D1%8C}
\bibitem{laine2010} Laine S., Karras T. Efficient sparse voxel octrees. Proceedings of the 2010 ACM SIGGRAPH Symposium on Interactive 3D Graphics and Games, 2010.
\bibitem{shirley2020} Shirley P. Ray Tracing in One Weekend. \url{https://raytracing.github.io/books/RayTracingInOneWeekend.html}
\bibitem{yagel1992} Yagel R., Cohen D., Kaufman A. Discrete ray tracing. IEEE Computer Graphics and Applications, 1992.
\end{thebibliography}

\end{document}
